\section{The $b$-field: the antisymmetric part of the shift}
\label{d4:sec:bfield}

The hypothesis places the $b$-field of Yang and Mills \cite{YangMills54} at the origin of the mass. The preceding sections built the
undivided energy and its differentiation without any connection. This section asks where a connection enters the interior dynamics of
Sections~15--17 of \cite{ES17}, and what it can do there. In brief: a constant shift of the interior Lagrangian splits into an
$\eta$-self-adjoint part, the binding term $\kappa G$, and an $\eta$-antisymmetric part $b\in\so{3,3}$, and only $b$ is a connection
(\cref{d4:thm:covariant}). If $b$ commutes with $G$, it splits the rest frequency into a Zeeman pattern, but it can be gauged away along any
worldline (\cref{d4:thm:lG,d4:prop:larmor}). If $b$ anticommutes with $G$, it shifts the frequency and always destroys the normal modes
(\cref{d4:thm:pG}), unless a commuting rotation screens it (\cref{d4:prop:gyro}). On its own, $b$ does not bind (\cref{d4:prop:alone}). Units are
$\hbar=c=1$ unless a symbol is shown. The parameter $\tau$ has dimension $T$; $\kappa=c/\ell$ (equation~(7) of \cite{ES17}) and every eigenvalue
of $b$ are rates, of dimension $T^{-1}$; and $\hbar\kappa=\hbar c/\ell$ is a quantum of interior energy. Every finite-dimensional statement was
checked by \texttt{verify/d4\_bfield.py} (\cref{d4:tab:claims}).

\paragraph{Notation.} For $a<b$, $R_{ab}\in\so{3,3}$ is defined by $R_{ab}e_b=\eta_{bb}e_a$ and $R_{ab}e_a=-\eta_{aa}e_b$; on vectors it equals
$\half\ad(e_ae_b)$. The generators of $\lG$ (Theorem~12.1 of \cite{ES17}) are $L_i=\half\ad(e_{j+3}e_{k+3}-e_je_k)$ and
$K_i=\ad\bigl(\half(e_ke_{j+3}-e_je_{k+3})\bigr)$ for cyclic $(i,j,k)$. On $z=\bB+i\bE$ they act as $z\mapsto\mathbf e_i\times z$ and
$z\mapsto i\,\mathbf e_i\times z$, so $K_i=GL_i$ and
\begin{equation}
\label{d4:eq:w}
b=\textstyle\sum_i\bigl(w^{\rm r}_iL_i+w^{\rm b}_iK_i\bigr)\in\lG\quad\text{acts as}\quad z\mapsto w\times z,\qquad w=w^{\rm r}+i\,w^{\rm b}\in\C^3 .
\end{equation}
In particular $R_{12}-R_{45}=-L_3$. The grading root is $\mathcal K=\half\ad\bigl(\sum_ie_ie_{i+3}\bigr)$, with $G=-\mathcal K\eta$
(Lemma~3.4 of \cite{ES17}). We write $\sigma_a$ for the boost exchanging $e_a$ and $e_{a+3}$, and $A^\dagger=\eta^{-1}A^{\!\top}\eta$. The
interior Lagrangian with a constant shift is
\begin{equation}
\label{d4:eq:L}
\mathcal L=\frac m2\,\eta(u,u),\qquad u=\dot X+SX,\qquad S=b+\kappa G,\qquad b\in\so{3,3},
\end{equation}
with $p=m\eta u$ and $\pi=\eta^{-1}p=mu$. For $b=0$ this is the split-form Lagrangian of Theorem~16.2 of \cite{ES17}.

\subsection{Two summands, and which one is a connection}
\label{d4:sec:split}

\begin{proposition}[Two summands]
\label{d4:prop:split}
$\End(V)=\so{3,3}\oplus\{A:A=A^\dagger\}$, of dimensions $15+21$. The part of $S=b+\kappa G$ in the first summand is $b$, and the part in the
second is $\kappa G$. The binding root lies in the second summand ($G^\dagger=G$) and the grading root in the first ($\mathcal K^\dagger=-\mathcal K$).
\end{proposition}

\begin{proof}
$A\mapsto A^\dagger$ is an involution, and $A$ lies in its $\mp1$ eigenspace exactly when $\eta A$ is antisymmetric (respectively symmetric). The
statements about the roots are Proposition~3.3 of \cite{ES17}.
\end{proof}

The decomposition itself is linear algebra valid for any nondegenerate form, so it is inherited. That $G$ and $\mathcal K$ fall on opposite sides
is fixed (Theorem~17.3 of \cite{ES17}). The next theorem gives the decomposition its dynamical meaning.

\begin{theorem}[Covariant factorisation]
\label{d4:thm:covariant}
Let $b\in\so{3,3}$ be constant and $D:=\frac{d}{d\tau}+b$.
\begin{enumerate}[label=(\roman*),leftmargin=2.2em]
\item The Euler--Lagrange equations of \eqref{d4:eq:L} are $(D-\kappa G)\pi=0$ with $\pi=m(D+\kappa G)X$, that is
\begin{equation}
\label{d4:eq:covEL}
(D-\kappa G)(D+\kappa G)X=0\qquad\Longleftrightarrow\qquad D^2X+\kappa^2X+\kappa\,(D_\tau G)\,X=0,\qquad D_\tau G:=[b,G].
\end{equation}
\item Let $X=g(\tau)Y$ with $g(\tau)\in O(3,3)$. Then $\bigl(\tfrac d{d\tau}+S\bigr)(gY)=g\bigl(\tfrac d{d\tau}+S^g\bigr)Y$, where $S^g=g^{-1}Sg+g^{-1}\dot g$. The
$\so{3,3}$ part $b^g=g^{-1}bg+g^{-1}\dot g$ transforms as a connection. The self-adjoint part $\kappa g^{-1}Gg$ transforms as a tensor. And
$(D_\tau G)^g=g^{-1}(D_\tau G)g$.
\end{enumerate}
\end{theorem}

\begin{proof}
(i) $\partial\mathcal L/\partial\dot X=p$ and $\partial\mathcal L/\partial X=S^{\!\top}p$, so $\dot p=S^{\!\top}p$. Equivalently
$\dot\pi=S^\dagger\pi=(-b+\kappa G)\pi$, which is $(D-\kappa G)\pi=0$. Moreover $\pi=m(D+\kappa G)X$, and since $G$ is constant,
$D(GX)-GDX=[b,G]X$ and $G^2=-\id$.

(ii) Differentiating $g^{\!\top}\eta g=\eta$ gives $(g^{-1}\dot g)^\dagger=-g^{-1}\dot g$, so the inhomogeneous term lies in $\so{3,3}$, and
$(g^{-1}Ag)^\dagger=g^{-1}A^\dagger g$ preserves both summands. Finally, $\frac d{d\tau}(g^{-1}Gg)+[b^g,g^{-1}Gg]=g^{-1}[b,G]g$. The check uses fifteen
symbolic coordinates for $b$ and a Cayley curve $g(\tau)=(\id-\tau A)^{-1}(\id+\tau A)$.
\end{proof}

The theorem is fixed. The factorisation of Theorem~16.2 of \cite{ES17} survives every connection once $d/d\tau$ is replaced by $D$. The only
new term is $\kappa(D_\tau G)X$, the same term as in the curved extension at the end of Section~15 of \cite{ES17}, with the connection now
internal. A Yang--Mills potential can therefore sit in the $\so{3,3}$ summand and nowhere else. No gauge transformation creates, removes or
rescales the binding part, since that part is a tensor. If the shift were the grading root, $\kappa\mathcal K$ would itself be a connection and
could be gauged away (\cref{d4:prop:alone}). The dictionary entry ``$b\leftrightarrow$ transport part of the shift'' is therefore correct,
and forced.

\begin{proposition}[Commuting and anticommuting parts]
\label{d4:prop:lp}
$\Ad_G:b\mapsto GbG^{-1}$ is an involution of $\so{3,3}$, so $\so{3,3}=\lG\oplus\mathfrak p_G$, where $\lG=\{[b,G]=0\}$ has dimension $6$ and
$\mathfrak p_G=\{bG=-Gb\}$ has dimension $9$. In the variable $z$,
\[
\lG=\{z\mapsto(P+iQ)z\}\cong\so{3,\C},\qquad \mathfrak p_G=\{z\mapsto(X+iY)\bar z\},
\]
with $P,Q,X$ real antisymmetric and $Y$ real symmetric. The Killing form $4\Tr$ has signature $(3,3)$ on $\lG$. On $\mathfrak p_G$ it is
negative on the three directions $Y=0$, which are compact counter-rotations of the two triples such as $R_{12}+R_{45}$, and positive on the
six boosts $X=0$. Both $\mathcal K$ ($z\mapsto-i\bar z$) and the pairing boosts $\sigma_a$ lie in $\mathfrak p_G$.
\end{proposition}

\begin{proof}
$G^{-1}=-G$ and $G^\dagger=G$ give $(GbG^{-1})^\dagger=-GbG^{-1}$ and $\Ad_G^2=\Ad_{-\id}=\id$. In $(\bB,\bE)$ blocks, commuting with $G$ means the form
$\bigl(\begin{smallmatrix}P&-Q\\Q&P\end{smallmatrix}\bigr)$, and anticommuting means $\bigl(\begin{smallmatrix}X&Y\\Y&-X\end{smallmatrix}\bigr)$; the
condition $b^\dagger=-b$ gives the symmetry types. The Killing form was computed in the adjoint representation.
\end{proof}

This is fixed. The pairing boosts that tear the binding at a threshold (Proposition~16.7 of \cite{ES17}) are anticommuting, and so is the
grading root.

\subsection{The phase-space generator and a criterion for boundedness}
\label{d4:sec:M}

\begin{lemma}[Bookkeeping]
\label{d4:lem:M}
On $(X,\pi)$ the flow is generated by $M=\bigl(\begin{smallmatrix}-S&\id/m\\0&S^\dagger\end{smallmatrix}\bigr)$, with $S^\dagger=-b+\kappa G$. It is the
Hamiltonian flow of the Noether energy $H(X,p)=\tfrac1{2m}p^{\!\top}\eta^{-1}p-p^{\!\top}SX$. $H$ vanishes on the invariant Lagrangian subspace $p=0$,
and $\det(\lambda-M)=\chi_S(\lambda)\chi_S(-\lambda)$, where $\chi_S(\lambda)=\det(\lambda-S)$.
\end{lemma}

\begin{proof}
$\dot X=\pi/m-SX$ and $\dot\pi=S^\dagger\pi$. $M$ is block triangular, and $S^\dagger$ is similar to $S^{\!\top}$.
\end{proof}

This is Hamiltonian bookkeeping, inherited. The classification rests on the following criterion.

\begin{theorem}[When the interior motion stays bounded]
\label{d4:thm:criterion}
Let $S$ be any constant real $6\times6$ matrix and $D=S^\dagger$.
\begin{enumerate}[label=(\roman*),leftmargin=2.2em]
\item $M$ is semisimple if and only if $S$ is semisimple and, for every $\mu\in\C$, the eigenspaces $E_\mu(D)$ and $E_{-\mu}(D)$ are orthogonal for
the bilinear extension of $\eta$.
\item Every solution is bounded on $\R$ if and only if, in addition, $\spec S\subset i\R$. Then (i) says that every eigenspace of $D$ is totally
isotropic for $h(u,y)=\bar u^{\!\top}\eta y$. Equivalently, every momentum normal mode $\pi=\Re(ye^{i\omega\tau})$, superpositions within a
degenerate frequency included, has period mean $\langle\eta(\pi,\pi)\rangle=\half h(y,y)=0$.
\end{enumerate}
\end{theorem}

\begin{proof}
$M$ is semisimple if and only if $A=-S$ (its restriction to $\pi=0$) and $D$ (the quotient map) are semisimple and $M$ is similar to
$\operatorname{diag}(A,D)$. By Roth's theorem \cite{Roth52}, the latter holds if and only if $AY-YD=-\id/m$ is solvable. Diagonalise
$A=P\Lambda_AP^{-1}$ and $D=Q\Lambda_DQ^{-1}$. The equation is then solvable if and only if $(P^{-1}Q)_{ij}=0$ whenever
$(\Lambda_A)_i=(\Lambda_D)_j=:\mu$. A row $\ell^{\!\top}$ of $P^{-1}$ satisfies $S^{\!\top}\ell=-\mu\ell$, and $S^{\!\top}=\eta S^\dagger\eta^{-1}$ puts
$w=\eta^{-1}\ell$ in $E_{-\mu}(D)$; the condition is $\ell^{\!\top}y=\eta(w,y)=0$ for $y\in E_\mu(D)$. For (ii), a linear flow is bounded on $\R$ exactly when
it is semisimple with imaginary spectrum, and then $E_{-\mu}(D)=\overline{E_\mu(D)}$.
\end{proof}

The theorem is fixed. Its necessity half is the split virial identity (Theorem~16.1 of \cite{ES17}: a bound orbit has null momentum on
average). The new content is that the averaged virial condition, imposed mode by mode, is also sufficient. For $S=\kappa G$ it reduces to the
isotropy statement of Proposition~16.3 of \cite{ES17}. On seven test connections we compared exact semisimplicity, exact Sylvester
solvability over $\Q$, and the isotropy condition evaluated numerically to $10^{-8}$; all three agreed.

\subsection{Connections commuting with $G$: a Zeeman pattern without a scale}
\label{d4:sec:lG}

\begin{theorem}[Connections commuting with $G$]
\label{d4:thm:lG}
Let $b\in\lG$ act as in \eqref{d4:eq:w}.
\begin{enumerate}[label=(\roman*),leftmargin=2.2em]
\item $Y=-G/(2\kappa m)$ solves $AY-YD=-\id/m$, so $M$ is semisimple if and only if $S$ is.
\item $\chi_S(\lambda)=(\lambda^2+\kappa^2)\bigl((\lambda-i\kappa)^2+w\cdot w\bigr)\bigl((\lambda+i\kappa)^2+\overline{w\cdot w}\bigr)$.
\item Every solution is bounded if and only if $w=0$ or $w\cdot w>0$. The frequencies are then $\kappa$ and $|\kappa\pm s|$, where
$s=\sqrt{w\cdot w}$. For $w\cdot w=0\neq w$ the zero-momentum flow grows like $\tau^2$; for $w\cdot w\notin[0,\infty)$ it grows exponentially.
\item For the rotation $b=\beta(R_{12}-R_{45})=-\beta L_3$, the spectrum of $M$ is $\pm i\kappa$, $\pm i(\kappa+\beta)$, $\pm i(\kappa-\beta)$, each
double. For the boost $b=\beta K_3$, the exponents are $\pm\beta\pm i\kappa$: complex rather than real, so the orbits spiral outward at the rest
frequency.
\end{enumerate}
\end{theorem}

\begin{proof}
(i) With $[b,G]=0$: $-(b+\kappa G)Y-Y(-b+\kappa G)=\frac1{2\kappa m}(bG-Gb+2\kappa G^2)=-\frac1m$.

(ii) On $(V,G)\cong\C^3$, $S$ is $z\mapsto i\kappa z+w\times z$. The identity $(w\times)^2v=w(w\cdot v)-(w\cdot w)v$ gives
$(w\times)^3=-(w\cdot w)(w\times)$, so $w\times$ has eigenvalues $0$ and $\pm is$ with $s^2=w\cdot w$. The real characteristic polynomial of a
complex-linear map is its complex characteristic polynomial times the conjugate one.

(iii) The exponents $i\kappa$ and $i(\kappa\pm s)$ are imaginary exactly when $s$ is real, and they are distinct when $s\neq0$. If
$w\cdot w=0\neq w$, then $N=w\times$ has $N^2\bar w=|w|^2w\neq0$, so $S=i\kappa+N$ is not semisimple on the invariant subspace $\pi=0$.

(iv) Here $s=\pm\beta$, respectively $s=\pm i\beta$. The characteristic polynomials were checked symbolically, and semisimplicity exactly at
$(\kappa,\beta)=(1,\tfrac37)$ and $(1,\tfrac32)$.
\end{proof}

The theorem is fixed. It lifts the degeneracy noted after Proposition~16.7 of \cite{ES17} (``the single frequency is degenerate across the
three pairs''): a rotation splits the rest frequency into the triplet $\kappa-\beta$, $\kappa$, $\kappa+\beta$ and keeps every orbit bounded. The
criterion in (iii) also has an interpretive reading. Through the isomorphism $(V,\beta_0)\cong\lG$ of Theorem~17.5 of \cite{ES17}
($L_i\leftrightarrow e_i$, $K_i\leftrightarrow e_{i+3}$), the connection is itself a frame vector, and $w\cdot w$ is its invariant $q$. A commuting
connection keeps the binding exactly when it is of magnetic type: $\bE_b\cdot\bB_b=0$ and $|\bB_b|>|\bE_b|$.

\begin{proposition}[The commuting part is a rotating frame]
\label{d4:prop:larmor}
Let $b\in\lG$. Then $\bigl(\tfrac d{d\tau}+b+\kappa G\bigr)(e^{-b\tau}Y)=e^{-b\tau}(\dot Y+\kappa GY)$. The exponents are therefore $\pm i\kappa$ shifted by the
eigenvalues of $-b$, and with $p_X=e^{b^{\!\top}\tau}p_Y$,
\[
H_{\rm lab}=H_{\rm rot}-p_Y^{\!\top}bY.
\]
\end{proposition}

\begin{proof}
$e^{-b\tau}$ is an isometry that commutes with $b$ and $G$. Moreover $e^{b\tau}\eta^{-1}e^{b^{\!\top}\tau}=\eta^{-1}$ and
$e^{b\tau}(b+\kappa G)e^{-b\tau}=b+\kappa G$.
\end{proof}

This is Larmor's theorem \cite{Larmor1897}, inherited: a rotating frame trades the field for an angular velocity. The term $Q_b=p_Y^{\!\top}bY$ is
the Noether charge of the $\lG$ rotation, so $H_{\rm lab}=H_{\rm rot}-Q_b$ is a Hamiltonian with a chemical potential for that charge. The
consequence is interpretive but sharp. The Zeeman pattern of \cref{d4:thm:lG} labels the rest-frequency levels by their $\lG$ charge. On one
worldline it has no gauge-invariant content; it becomes observable only against a reference that $b$ does not rotate, just as the Zeeman
effect is observed against the laboratory. A commuting $b$ supplies a spectrum but no new invariant scale, and $\beta/\kappa$ is free.

\subsection{Connections anticommuting with $G$: tearing, and gyroscopic screening}
\label{d4:sec:pG}

\begin{theorem}[Connections anticommuting with $G$]
\label{d4:thm:pG}
Let $b\in\mathfrak p_G$ with $b\neq0$.
\begin{enumerate}[label=(\roman*),leftmargin=2.2em]
\item $(\pm b+\kappa G)^2=b^2-\kappa^2$, and
$M^2=\operatorname{diag}(b^2-\kappa^2,\,b^2-\kappa^2)+N$ with $N=\bigl(\begin{smallmatrix}0&-2b/m\\0&0\end{smallmatrix}\bigr)\neq0$. Hence $M$ is never
semisimple, and some solutions grow at least linearly.
\item For the compact element $b=\beta(R_{12}+R_{45})$, $\det(\lambda-M)=(\lambda^2+\kappa^2)^2(\lambda^2+\kappa^2+\beta^2)^4$. The exponents
$\pm i\sqrt{\kappa^2+\beta^2}$ have algebraic multiplicity four and geometric multiplicity two, and $\pm i\kappa$ are semisimple. The bounded
initial data are exactly those with $\pi_1=\pi_2=\pi_4=\pi_5=0$, that is, with no momentum on the support of $b$. The zero-momentum orbits
are bounded, and on the support of $b$ they turn at $\sqrt{\kappa^2+\beta^2}$.
\item For a pairing-plane boost $b=c\,\sigma_a$, the exponents in that plane are $\pm\sqrt{c^2-\kappa^2}$ (Proposition~16.7 of \cite{ES17}), with
Jordan blocks for every $c\neq0$. This holds below the threshold (at $c=\tfrac12$, $\kappa=m=1$, $\rank(M^2+\tfrac34)=2$ on the plane), above it
($c=2$), and at it, where $S$ restricted to the plane is nilpotent.
\item The grading root read as a connection satisfies $(c\mathcal K+\kappa G)^2=(c^2-\kappa^2)\id$ on all of $V$.
\end{enumerate}
\end{theorem}

\begin{proof}
(i) $\{b,G\}=0$ gives both squares, and $M^2$ has off-diagonal block $(A+D)/m=-2b/m$. Write $b^2-\kappa^2=Q_s+Q_n$ (Jordan--Chevalley). Then
$M^2=\operatorname{diag}(Q_s,Q_s)+\bigl(\operatorname{diag}(Q_n,Q_n)+N\bigr)$, and the summands commute because $Q_n$ is a polynomial in $b$. By uniqueness,
$M^2$ is semisimple only if $\operatorname{diag}(Q_n,Q_n)+N=0$, which is impossible because $N\neq0$ is off-diagonal. If $M$ were semisimple, so would
be $M^2$.

(ii)--(iv) are exact rank and null-space computations over $\Q(i)$ at $(\kappa,\beta)=(4,3)$, where the shifted frequency is $5$, together with the
symbolic characteristic polynomials. Part (i) was checked with nine symbolic coordinates on $\mathfrak p_G$.
\end{proof}

Part (i) is an obstruction: no anticommuting connection is compatible with bounded normal modes. Parts (ii)--(iv) are fixed. The
anticommuting part does change the rate, to $\sqrt{\kappa^2+\beta^2}$ for a compact element and $\sqrt{\kappa^2-c^2}$ for a boost. In the momentum
sector, however, it does what the positive form does in Theorem~15.5 of \cite{ES17}: imaginary exponents carried by Jordan blocks.

\begin{proposition}[What a gauge transformation removes]
\label{d4:prop:remove}
Let $b=b_l+b_p\in\lG\oplus\mathfrak p_G$ be constant. Some $g(\tau)\in O(3,3)$ carries $(b,G)$ to $(0,G')$ with $G'$ constant if and only if $[b,G]=0$. For
$b\in\mathfrak p_G$, $e^{b\tau}Ge^{-b\tau}=e^{2b\tau}G$. The gauge-covariant remainder is $D_\tau G=[b,G]=2b_pG$: along a worldline, $b_l$ is pure gauge
and $b_p$ is not.
\end{proposition}

\begin{proof}
$b^g=0$ forces $g=e^{-b\tau}g_0$, and then $G'=g_0^{-1}e^{b\tau}Ge^{-b\tau}g_0$ has derivative $g_0^{-1}e^{b\tau}[b,G]e^{-b\tau}g_0$. If $b$ anticommutes
with $G$, then $Gf(b)=f(-b)G$ for every power series $f$.
\end{proof}

This is fixed. A connection on a line has no curvature, and its only invariants are holonomies around closed loops (inherited). The
non-abelian curvature $[b_\mu,b_\nu]$ of constant potentials needs two directions, while a worldline sees only $b_\tau=\dot x^\mu b_\mu$. Once
the binding root is present, the invariant is $D_\tau G$, which is exactly the anticommuting part.

One might expect boundedness to force $b\in\lG$. It does not.

\begin{proposition}[Gyroscopic screening]
\label{d4:prop:gyro}
Let $b=aL_3+b_Y$, where $b_Y\in\mathfrak p_G$ is the boost $z\mapsto iY\bar z$ with $Y=\bigl(\begin{smallmatrix}c&d&0\\d&-c&0\\0&0&0\end{smallmatrix}\bigr)$,
and let $r^2=c^2+d^2>0$. Then $[b,G]\neq0$ and
\[
\chi_S(\lambda)=(\lambda^2+\kappa^2)\bigl(\lambda^2+(\kappa-a)^2-r^2\bigr)\bigl(\lambda^2+(\kappa+a)^2-r^2\bigr).
\]
If $a\neq0$ and $r<|\kappa\pm a|$, every solution is bounded; for $a=0$ some solutions are unbounded (\cref{d4:thm:pG}). Numerically, the
set of $b\in\so{3,3}$ for which all motion is bounded is locally a $10$-dimensional submanifold of $\so{3,3}$.
\end{proposition}

\begin{proof}
On $U=\langle e_1,e_2,e_4,e_5\rangle$, $K_3=GL_3$ is an involution commuting with $G$, $L_3$ and $b_Y$. On its eigenplanes $U_\pm$ we have
$L_3=\mp G$, so $S|_{U_\pm}=(\kappa\mp a)G+b_Y$ and $S^\dagger|_{U_\pm}=(\kappa\pm a)G-b_Y$. Each of these squares to a negative multiple of the identity,
because $b_Y^2=r^2$ on $U$ and $b_Y$ anticommutes with $G$, so each is semisimple with imaginary spectrum. Because $K_3$ is $\eta$-antisymmetric with
$K_3^2=\id_U$, the planes $U_\pm$ are totally $\eta$-null. For $a\neq0$ they carry different frequencies, so every eigenvector of $S^\dagger$ lies in
$U_+$, in $U_-$, or in the $b$-free plane $\langle e_3,e_6\rangle$, and is $h$-isotropic. \Cref{d4:thm:criterion} gives boundedness, which was also
checked exactly at $(a,c,d)=(\tfrac12,\tfrac14,0)$ and $(\tfrac13,\tfrac15,\tfrac17)$ with $\kappa=1$. For the dimension: for simple frequencies,
\cref{d4:thm:criterion}(ii) imposes five independent real conditions (two for imaginary spectrum, after the trace relation, and three for
isotropy). Their Jacobian has rank five at four solutions found by least squares with $|b_p|=0.3$ (relative singular-value threshold
$10^{-4}$), and these solutions stay bounded up to $\tau=10^3$.
\end{proof}

The family is fixed; the dimension $10$ is numerical. The mechanism is the classical gyroscopic stabilisation of a saddle
\cite{ThomsonTait1879,Kirillov2013}. The boost acts inside the two light-like planes of $K_3$, and the rotation detunes their frequencies, so the
virial $\eta(\pi,\pi)=2\eta(\pi_+,\pi_-)$ averages to zero. This corrects the working dictionary in two ways. An anticommuting part need not cause
secular growth when a rotation detunes the planes it couples. And boundedness is a codimension-five condition that does not single out $\lG$;
what singles out $\lG$ is gauge removability (\cref{d4:prop:remove}). The screened frequencies $\sqrt{(\kappa\mp a)^2-r^2}$ still require
$\kappa$ or $a$ to be nonzero, so screening supplies no scale either.

\begin{proposition}[A connection alone does not bind]
\label{d4:prop:alone}
\begin{enumerate}[label=(\roman*),leftmargin=2.2em]
\item Let $\kappa=0$ and let $b\in\so{3,3}$ be any constant connection, $b=c\mathcal K$ included. Then $X=e^{-b\tau}Y$ gives
$\mathcal L=\tfrac m2\eta(\dot Y,\dot Y)$ and $\ddot Y=0$. The generator $M=\bigl(\begin{smallmatrix}-b&\id/m\\0&-b\end{smallmatrix}\bigr)$ has
$\det(\lambda-M)=\chi_{-b}(\lambda)^2$ and is never semisimple: the motion is free drift, and there is no rest frequency.
\item The rescaling $(\tau,\kappa,b,m)\mapsto(s\tau,\kappa/s,b/s,sm)$ maps $M$ to $M/s$ and solutions to solutions, and it leaves $\eta$, $G$, $\lG$,
$\mathfrak p_G$ and $\so4_\Delta$ unchanged. No value of $\kappa$, hence of $\ell=c/\kappa$, is selected, with or without a constant $b$.
\end{enumerate}
\end{proposition}

\begin{proof}
(i) $\dot X+bX=e^{-b\tau}\dot Y$. Write $-\operatorname{diag}(b,b)=S'+N'$ (Jordan--Chevalley). Then $M=S'+(N'+N)$ with
$N=\bigl(\begin{smallmatrix}0&\id/m\\0&0\end{smallmatrix}\bigr)$, and $N'+N\neq0$ because $N$ is off-diagonal and $N'$ is not. (ii) is direct.
\end{proof}

Part (i) is inherited, since every connection on a line is pure gauge. Its consequence is an obstruction: a Yang--Mills potential on a
worldline, with no self-adjoint shift, produces free motion. Part (ii) is an obstruction for classical constant fields. For the $b$-field it
restates what Section~17 of \cite{ES17} finds for the closure: no ratio of scales is fixed by the Clifford structure.

\subsection{The split virial identity and the vacuum}
\label{d4:sec:vacuum}

\begin{proposition}[Split virial with a connection]
\label{d4:prop:virial}
For every constant $S$, and in particular for every $b\in\so{3,3}$, $\frac d{d\tau}\eta(X,\pi)=\eta(\pi,\pi)/m$ along every solution.
\end{proposition}

\begin{proof}
$\eta(\pi/m-SX,\pi)+\eta(X,S^\dagger\pi)=\eta(\pi,\pi)/m$. This was checked symbolically with $36$ coordinates for $S$, and in integral form along
numerical solutions (DOP853, tolerance $10^{-12}$, error below $10^{-8}$) for $b$ in $\lG$, in $\mathfrak p_G$, and of mixed type.
\end{proof}

This re-verifies Theorem~16.1 of \cite{ES17}, whose metric-compatible connection is our $b$, and it is fixed.

\begin{theorem}[Krein signatures and the zero-point energy]
\label{d4:thm:krein}
Suppose every solution is bounded, and let $\omega>0$ have multiplicity $k_\omega$ in $\spec S$. On the eigenspace $E_{i\omega}(M)$, the energy form
$v\mapsto\bar v^{\!\top}\mathcal Hv$, where $\mathcal H$ is the Hessian of $H$, has signature $(k_\omega,k_\omega)$. Hence, with positive-norm
quantisation of the normal modes, the zero-point energy $\sum_j(\pm)\half\hbar\omega_j$ cancels frequency by frequency, and every level has a lower
one.
\end{theorem}

\begin{proof}
On $E_{i\omega}$, $\mathcal Hv=-i\omega\mathbb Jv$, so the energy form equals $\omega$ times the Krein form $-i\bar v^{\!\top}\mathbb Jv$. The Krein form is
nondegenerate on $E_{i\omega}$, because eigenspaces of distinct imaginary eigenvalues are orthogonal for it \cite{YakubovichStarzhinskii75}. The
zero-momentum part of $E_{i\omega}$ has dimension equal to the multiplicity of $i\omega$ in $-S$, which is $k_\omega$, and $H$ vanishes on it
(\cref{d4:lem:M}). Since $M$ is semisimple, $\dim E_{i\omega}=2k_\omega$ by \cref{d4:lem:M}. A nondegenerate Hermitian form with a half-dimensional
isotropic subspace has signature $(k,k)$, and a mode of energy sign $\pm$ quantises to $\pm\hbar\omega(n+\half)$.
\end{proof}

The theorem is fixed. It was checked numerically (to $10^{-9}$) on the Zeeman triplet at $\beta=\tfrac37$ and at $\beta=\tfrac32$, on the elliptic
complex $w=(2,i,0)$, on the gyroscopic family, and at three constituents. It extends Theorem~18.6 of \cite{ES17}, where the zero-point energy on
$\R^{3,3}$ vanishes at $b=0$, to every bounded connection. Neutral-signature cancellation survives the splitting line by line, and so does the
absence of a ground state. The cancellation is forced by the perfect-square form \eqref{d4:eq:L}, because the zero-momentum half of each frequency
carries no energy. On the coordinate subfields of that theorem, where the zero-point energy is $\pm\tfrac32\hbar\kappa$ or $\pm\half\hbar\kappa$, $G$ maps
the subfield out of itself, and the theorem does not apply.

\begin{corollary}[No Fock vacuum with Jordan blocks]
\label{d4:cor:nofock}
Suppose $M$ is not semisimple at a frequency $\omega\neq0$, as for every nonzero $b\in\mathfrak p_G$. Then no linear symplectic change of
variables writes the Hamiltonian on that block as a sum of signed oscillators. Hence there is no Fock vacuum annihilated by normal-mode
lowering operators.
\end{corollary}

\begin{proof}
A sum of signed oscillators generates a semisimple flow, and symplectic changes of variables act by similarity.
\end{proof}

This is an obstruction. A model block with the same Jordan structure (checked) is $H=\omega(q_1p_2-q_2p_1)+\half(q_1^2+q_2^2)$, whose two terms
Poisson-commute \cite{Williamson36}. Its quantisation, a rotation generator plus a multiplication operator, has continuous spectrum; this
last statement is inherited and was not re-derived.

\subsection{Three constituents: the Yang--Mills-type algebra compatible with binding}
\label{d4:sec:three}

We use the generators $E_{6(k-1)+a}=\omega^{\otimes(k-1)}\otimes e_a\otimes\id^{\otimes(3-k)}$ of equation~(2) of \cite{ES17} on $\R^{512}$, with the
adjoint action of bivectors on the $18$ vectors computed exactly, and set $G_3=G\oplus G\oplus G$. We follow the numbering of Section~17 of
\cite{ES17}: $W=\{7,13,14,15\}$, the colour directions are $\{10,11,12,16,17,18\}$, $E=\operatorname{span}\{7,8,9,13,14,15\}$, and
$\mathrm{rot}(i,j)=e_{i'}e_{j'}-e_ie_j$ acts by $\half\ad$.

\begin{proposition}[The binding-compatible algebra]
\label{d4:prop:PS}
The $E_A$ satisfy the relations of $\Cl(9,9)$. $G_3$ is an anti-isometry with $G_3^2=-\id$, and $\dim\bigl(\Cent(G_3)\cap\so{9,9}\bigr)=72$. For
four-dimensional $U\subset\operatorname{span}\{2,3,7,8,9,13,14,15\}$, neither $\so U$ nor $\so6_C$ meets $\Cent(G_3)$, while their sum meets it in
dimension $\binom d2$, $d=\dim(U\cap E)$ (Theorem~17.21 of \cite{ES17}). We checked the three cases $U=W$ (dimension $6$), $U=\{3,7,8,9\}$
($3$) and $U=\{2,3,7,8\}$ ($1$). For $U=W$ the intersection is
$\so4_\Delta=\operatorname{span}\{\mathrm{rot}(i,j):i<j\in W\}$. With $(a,b,c,d)=(7,13,14,15)$, the elements
$T^\pm_1=\half(\mathrm{rot}(b,c)\pm\mathrm{rot}(a,d))$, $T^\pm_2=\half(\mathrm{rot}(a,c)\mp\mathrm{rot}(b,d))$ and $T^\pm_3=\half(\mathrm{rot}(a,b)\pm\mathrm{rot}(c,d))$
satisfy $[T^+_i,T^-_j]=0$ and $[T^\pm_1,T^\pm_2]=-T^\pm_3$ cyclically. The trace form is negative definite, so $\so4_\Delta=\su2_+\oplus\su2_-$ is
compact. Each $T^\pm_3$ has eigenvalues $\pm\tfrac i2$, four times each, and $0$, ten times.
\end{proposition}

This re-verifies Theorem~17.21 of \cite{ES17} on the eighteen vectors, and it is fixed. Each $\su2$ factor locks a weak rotation of $W$ to the
colour rotation of the partners $\{10,16,17,18\}$. Reading $\so4_\Delta$ as ``Yang and Mills' isospin, doubled'' is an interpretation of
$\so4\cong\su2\oplus\su2$.

\begin{proposition}[Spectra on the eighteen-dimensional frame]
\label{d4:prop:so4spec}
Let $\kappa=5$ and $m=1$.
\begin{enumerate}[label=(\roman*),leftmargin=2.2em]
\item For $b=\beta T^+_3$ with $\beta=2$, $\chi_S=(\lambda^2+25)^5(\lambda^2+16)^2(\lambda^2+36)^2$. The frequency is $\kappa$ on the five complex
singlet directions and $\kappa\mp\beta/2$ on the two $\su2_+$ doublets, which have weights $\pm\half$.
\item For $b=\beta_+T^+_3+\beta_-T^-_3$ with $(\beta_+,\beta_-)=(2,4)$, the lines are $\kappa\pm\half(\beta_+\pm\beta_-)=8,6,4,2$, together with $\kappa$.
\item In both cases $Y=-G_3/(2\kappa m)$ solves the Sylvester equation, $M$ is semisimple, the Krein signatures are $(k,k)$, and the zero-point
energy vanishes.
\item A weak-only rotation $R_{7,13}$ or a colour-only rotation $R_{10,16}$ gives a non-semisimple $M$, hence secular growth.
\end{enumerate}
\end{proposition}

This is fixed: the characteristic polynomials and semisimplicity tests are exact on $36\times36$ integer matrices, and the Krein signatures are
numerical. Every constant $\so4_\Delta$-valued connection is compact and commutes with $G_3$, so by the argument of \cref{d4:thm:lG}(i) it keeps the
interior bounded. The splitting follows the $\su2$ weights, $\beta/\kappa$ is free, and by \cref{d4:prop:larmor} the whole pattern can be
removed along one worldline. A Pati--Salam connection in one factor alone tears the binding. Only the locked diagonal is compatible with it,
which is \cref{d4:prop:remove} at three constituents.

\subsection{Quantum Yang--Mills facts, and the channels of the paper's vacuum}
\label{d4:sec:qym}

The only place left for $b$ to generate a scale is the quantum theory. We re-verify the facts used by the earlier session.

\begin{proposition}[Inherited facts]
\label{d4:prop:qym}
\begin{enumerate}[label=(\roman*),leftmargin=2.2em]
\item For $\su4$ with $\tr(T^AT^B)=\half\delta^{AB}$, $\sum_AT^AXT^A=\half\bigl(\tr(X)\id-\tfrac14X\bigr)$. This is $\tfrac{15}8=C_F$ on $\id$ and $-\tfrac18$ on
traceless $X$.
\item The channel factors $T_1\cdot T_2=\half[C_2(\text{pair})-C_2(R_1)-C_2(R_2)]$ are $+\tfrac38$ and $-\tfrac58$ on $\mathbf{10}$ and $\mathbf6$ in
$\mathbf4\otimes\mathbf4$; $-\tfrac{15}8$ and $+\tfrac18$ on $\mathbf1$ and $\mathbf{15}$ in $\bar{\mathbf4}\otimes\mathbf4$; and $+\tfrac14$ and $-\tfrac34$ on $\mathbf3$ and
$\mathbf1$ in $\mathbf2\otimes\mathbf2$ of $\su2$. Here $C_2=\tfrac{15}8,\tfrac92,\tfrac52,4$ on $\mathbf4,\mathbf{10},\mathbf6,\mathbf{15}$.
\item $\int_{k^2\le L^2}\frac{d^4k}{(2\pi)^4}\frac1{k^2+M^2}=\frac1{16\pi^2}\bigl[L^2-M^2\ln(1+L^2/M^2)\bigr]$ (Euclidean, cutoff $L$). The function
$f(x)=1-x^2\ln(1+x^{-2})$ decreases strictly from $1$ to $0$, so $1=\lambda f(x)$ has a root $x=M/L>0$ if and only if $\lambda>1$. At
$\lambda=\tfrac74$ the root is $x=0.5330234964\ldots$, which confirms the earlier value $0.533$.
\item The ladder critical coupling $\pi/(3C_F)$ \cite{FukudaKugo76,Miransky93} equals $8\pi/45\approx0.5585$ at $C_F=\tfrac{15}8$.
\item For $SU(2)$ with $N_f$ Dirac doublets, $\beta_0=(22-2N_f)/3$, which is $6$ at $N_f=2$ \cite{GrossWilczek73,Politzer73}, and
$\Lambda=\mu e^{-8\pi^2/(\beta_0g^2)}$ is independent of $\mu$.
\item The breaking $SU(4)\to Sp(4)$ ($N_f=2$) leaves $15-10=5$ Goldstone directions \cite{Peskin80}.
\end{enumerate}
\end{proposition}

\begin{proof}
(i) and (ii) were computed exactly with the generalised Gell-Mann basis and the Casimirs on $\mathrm{Sym}^2$, $\Lambda^2$ and
$\bar{\mathbf4}\otimes\mathbf4$. (iii) uses the measure $2\pi^2k^3dk$; monotonicity follows from $f'(x)=2x[(1+x^2)^{-1}-\ln(1+x^{-2})]<0$, since
$\ln(1+u)>u/(1+u)$ for $u>0$, and the root was found with mpmath to $25$ digits. (v) Along $dg^2/d\ln\mu=-2\beta_0g^4/(16\pi^2)$, the derivative
$\frac d{d\ln\mu}[\ln\mu-8\pi^2/(\beta_0g^2)]$ vanishes. (vi) is the dimension of $\{X\in\su4:X^{\!\top}\Omega+\Omega X=0\}$.
\end{proof}

All of this is inherited \cite{NambuJonaLasinio61,Slansky81}.

\begin{proposition}[Channels of the Pati--Salam bilinears]
\label{d4:prop:channels}
The Dirac bilinear $\overline{(\mathbf4,\mathbf2,\mathbf1)}\otimes(\mathbf4,\mathbf1,\mathbf2)=(\mathbf1,\mathbf2,\mathbf2)\oplus(\mathbf{15},\mathbf2,\mathbf2)$ has $\su4$
factors $-\tfrac{15}8$ and $+\tfrac18$. A Lorentz-scalar Majorana bilinear of identical Weyl fields $(\mathbf4,\mathbf1,\mathbf2)$ is symmetric in its
internal indices, so it lies in $\mathrm{Sym}^2(\C^4\otimes\C^2)=(\mathbf{10},\mathbf1,\mathbf3)\oplus(\mathbf6,\mathbf1,\mathbf1)$ (dimensions $30+6$), with
factors $(+\tfrac38,+\tfrac14)$ and $(-\tfrac58,-\tfrac34)$ respectively.
\end{proposition}

These representation statements are inherited; they were checked through joint Casimir eigenspaces on $\mathrm{Sym}^2(\C^8)$. Their reading is
interpretive. Under the most-attractive-channel heuristic of one-boson exchange \cite{RabyDimopoulosSusskind80}, where a negative factor means
attraction, the exchange favours the Dirac singlet $(\mathbf1,\mathbf2,\mathbf2)$ and the colour-breaking $(\mathbf6,\mathbf1,\mathbf1)$. It disfavours
$(\mathbf{15},\mathbf2,\mathbf2)$, which the Georgi--Jarlskog factor $-3$ needs (Proposition~21.2 of \cite{ES17}; \cite{GeorgiJarlskog79}), and it disfavours
$(\mathbf{10},\mathbf1,\mathbf3)$, which carries the paper's Standard-Model vacuum (Section~20 of \cite{ES17}). A vacuum generated dynamically by the
gauge interaction would, on this heuristic, point away from the paper's.

\subsection{Synthesis: where the $b$-field enters}
\label{d4:sec:synthesis}

\begin{enumerate}[label=(\alph*),leftmargin=2.2em]
\item \emph{The binding frequency comes from the self-adjoint part, never from a connection} (fixed; obstructed for $b$ alone). Of the two
parts of the shift, only $b$ is a gauge potential (\cref{d4:thm:covariant}), and on its own it produces free drift (\cref{d4:prop:alone}). The
rate $\kappa$, with $(\kappa G)^2=-c^2/\ell^2$, sits in the tensor part, which no gauge transformation creates or removes.
\item \emph{A $b$-field commuting with $G$ splits the rest frequency without supplying a scale} (fixed; the Zeeman reading is interpretive). It
gives $\kappa$ and $\kappa\pm\sqrt{w\cdot w}$, and the motion stays bounded exactly for magnetic-type connections (\cref{d4:thm:lG}). At three
constituents it gives the $\su2$-weight pattern of $\so4_\Delta$, the Yang--Mills-type algebra compatible with binding
(\cref{d4:prop:so4spec}). The zero-point cancellation survives line by line (\cref{d4:thm:krein}). But the splitting is a rotating-frame
effect, a chemical potential for the $\lG$ charge that can be removed along any worldline (\cref{d4:prop:larmor}), and $\beta/\kappa$ is free.
\item \emph{A $b$-field anticommuting with $G$ shifts the frequency but tears the binding} (obstructed; spectra fixed). It gives
$\sqrt{\kappa^2+\beta^2}$ for a compact element and $\sqrt{\kappa^2-c^2}$ for a boost, always with Jordan blocks, so there is no Fock vacuum
(\cref{d4:thm:pG,d4:cor:nofock}). A commuting rotation can screen such a part (\cref{d4:prop:gyro}); the screened frequencies still need $\kappa$
or $a$.
\item \emph{The only role left for $b$ in generating $m$ is quantum} (open). The arena leaves $\ell$ free (\cref{d4:prop:alone}), and classical
four-dimensional Yang--Mills theory has no dimensionful parameter. A scale arises only by dimensional transmutation,
$\Lambda_b=\mu\,e^{-8\pi^2/(\beta_0g^2)}$ \cite{ColemanWeinberg73,GrossWilczek73}.
\end{enumerate}

\begin{conjecture}[A transmuted binding length]
\label{d4:conj:transmute}
$\ell=\hbar c/\Lambda_b$ for the scale $\Lambda_b$ of a Yang--Mills $b$-field, so that $\hbar\kappa=\Lambda_b$ and $m=\Lambda_b/c^2$ up to a pure number.
\end{conjecture}

The conjecture is open: nothing in the arena derives it. A mechanism would need the rate $\kappa$ to be dynamical, as the rotor rate of
Section~17 of \cite{ES17} is, and coupled to a quantum gauge sector with a fixed $\beta_0$ and matching coupling. Even then, by
\cref{d4:prop:channels}, the most attractive channels of one-boson exchange are not those of the paper's vacuum.

Corrected so that it is true, the user's sentence reads as follows. In $E=mc^2$ the symmetric part is the self-adjoint binding shift $\kappa G$,
with square $-c^2/\ell^2$. The Yang--Mills $b$-field is the antisymmetric, transport part of the same shift. Classically it does not generate $m$:
when it commutes with $G$ it splits $mc^2$ into charge-labelled lines, and when it does not, it tears the binding. Whether its quantum scale is
the $\ell$ that the arena leaves free is open.

\begin{table}[t]
\centering
\small
\begin{tabular}{@{}>{\raggedright\arraybackslash}p{8.4cm}>{\raggedright\arraybackslash}p{3.6cm}>{\raggedright\arraybackslash}p{2.6cm}@{}}
\toprule
statement & class & checks\\
\midrule
\cref{d4:prop:split}: decomposition / roots on opposite sides & inherited / fixed & d4-0, d4-1\\
\cref{d4:thm:covariant}: covariant factorisation; $b$ connection, $\kappa G$ tensor & fixed & d4-1, d4-2\\
\cref{d4:prop:lp}: $\lG\oplus\mathfrak p_G$, signatures & fixed & d4-3\\
\cref{d4:lem:M}: generator and energy & inherited & d4-4\\
\cref{d4:thm:criterion}: Roth--virial criterion & fixed & d4-11\\
\cref{d4:thm:lG}: Zeeman triplet, magnetic-type criterion, boost spirals & fixed & d4-5, d4-6, d4-7\\
\cref{d4:prop:larmor}: rotating frame; Zeeman reading & inherited; interpretive & d4-13\\
\cref{d4:thm:pG}: (i) never semisimple / (ii)--(iv) spectra & obstructed / fixed & d4-8, d4-9, d4-10\\
\cref{d4:prop:remove}: removable part is $\lG$, invariant $D_\tau G$ & fixed & d4-13\\
\cref{d4:prop:gyro}: gyroscopic screening; bounded set of dimension $10$ & fixed (dimension numerical) & d4-12\\
\cref{d4:prop:alone}: no binding and no scale from $b$ & obstructed & d4-13, d4-25\\
\cref{d4:prop:virial}, \cref{d4:thm:krein}: virial; zero point cancels per line & fixed & d4-14, d4-15\\
\cref{d4:cor:nofock}: no Fock vacuum with Jordan blocks & obstructed & d4-16\\
\cref{d4:prop:PS}, \cref{d4:prop:so4spec}: three constituents & fixed & d4-17, d4-18, d4-19\\
\cref{d4:prop:qym}, \cref{d4:prop:channels}: Yang--Mills facts & inherited & d4-20 to d4-24\\
channel reading; isospin reading of $\so4_\Delta$ & interpretive & d4-21, d4-18\\
\cref{d4:conj:transmute}: $\ell=\hbar c/\Lambda_b$ & open & ---\\
\bottomrule
\end{tabular}
\caption{Claims of \cref{d4:sec:bfield}, their classes, and the checks printed by \texttt{d4\_bfield.py}.}
\label{d4:tab:claims}
\end{table}
